
pour tout @@M0@@ un @@M1@@-Module localement libre de type fini @@M2@@ en posant~:
@@M3@@
où @@M4@@ est le morphisme structural @@M5@@.

On voit alors qu'il suffit de montrer la proposition dans le cas où @@M6@@ est un fibré projectif sur @@M7@@, le cas d'un fibré de drapeaux s'en déduisant par récurrence. Or ce cas résulte immédiatement de 1.1 (resp. 2.1).

\textbf{Théorème 3.2.} \uline{Soit} @@M8@@ \uline{un schéma quasi-compact}. \uline{Le pré}-@@M9@@-\uline{anneau} @@M10@@, \uline{défini à partir des} @@M11@@\uline{-Modules localement libres de type fini}, \uline{est un} @@M12@@-\uline{anneau}~(V 2.4).

L'homomorphisme @@M13@@ étant additif, il suffit, pour vérifier que c'est un @@M14@@-homomorphisme, de faire la vérification pour des éléments du type @@M15@@, où @@M16@@ est un @@M17@@-Module localement libre de type fini. On est donc ramené à montrer que pour tout couple d'éléments @@M18@@ de @@M19@@, classes de @@M20@@-Modules localement libres de type fini @@M21@@, les relations V 2.4.2 sont vérifiées :
@@M22@@
@@M23@@

Comme @@M24@@ est un quasi-compact, on peut trouver un recouvrement fini de @@M25@@ par des ouverts disjoints tel que sur chacun @@M26@@ et @@M27@@ soient de rang constant. Or le @@M28@@ d'une somme directe finie d'ouverts est le produit des @@M29@@ de ces ouverts. Il suffit donc de vérifier les relations V 2.4.2 dans chacun des @@M30@@ des ouverts du recouvrement. On peut donc supposer @@M31@@ et @@M32@@ de rang constant sur @@M33@@.

